\documentclass[leqno,10pt,a4paper]{amsart}

%-------------------------------------------------------------------------------

\usepackage[full]{textcomp}
\usepackage{newpxtext}
\usepackage{times}
\usepackage{cabin} % sans serif
\usepackage[varqu,varl]{inconsolata} % sans serif typewriter
\usepackage[bigdelims,vvarbb]{newpxmath} % bb from STIX
\usepackage[cal=cm,bb=esstix,bbscaled=1.05]{mathalfa} % mathcal
%\usepackage[margin=2.54cm]{geometry}
\usepackage[bottom=1.25in,top=1.1in,left=1.25in,right=1.25in]{geometry}
\setlength{\footskip}{20pt}
\usepackage{savesym}
\let\savedegree\bigtimes
\let\bigtimes\relax
\usepackage{mathabx}
\let\bigtimes\savedegree
\usepackage[usenames,dvipsnames]{color}
\usepackage{hyperref}
\hypersetup{
 colorlinks,
 linkcolor={blue!90!black},
 citecolor={red!80!black},
 urlcolor={blue!50!black}
}
\usepackage{tikz}
\usetikzlibrary{decorations.markings}
\usepackage{mathrsfs}
\usepackage[all,cmtip]{xy}
\usepackage{mleftright}
\usepackage{booktabs}
\usepackage{cite}
\usepackage{enumitem}


%------------------------------------------------------------------------------
% added by me

\usepackage{graphicx}
\usepackage{subfigure}

\usepackage[outdir=./]{epstopdf}
\usepackage{caption}
\captionsetup{
font=small,
justification   = raggedright}

\usepackage{placeins}
\usepackage{comment}




%-------------------------------------------------------------------------------
\renewcommand{\baselinestretch}{1.1}
\setlist[enumerate]{labelsep=*, leftmargin=1.5pc}
\setlist[enumerate]{label=\normalfont(\roman*), ref=\roman*}
%-------------------------------------------------------------------------------

\newtheorem{thm}{Theorem}[section]
\newtheorem{lem}[thm]{Lemma}
\newtheorem{cor}[thm]{Corollary}
\newtheorem{prop}[thm]{Proposition}
%-------------------------------------------------------------------------------
\theoremstyle{definition}
\newtheorem{example}[thm]{Example}
\newtheorem{remark}[thm]{Remark}
\newtheorem{definition}[thm]{Definition}
\newtheorem{conjecture}[thm]{Conjecture}
\newtheorem{question}[thm]{Question}
\newtheorem{problem}[thm]{Problem}
\numberwithin{equation}{section}
%-------------------------------------------------------------------------------
\newcommand{\mypsdraft}{\psdraft}
\renewcommand{\mypsdraft}{\psfull}
\newcommand{\mypsfull}{\psfull}


\newcommand{\IGNORE}[1]{}
\newcommand{\ignore}[1]{}
\newcommand{\veps}{\varepsilon}
\newcommand{\tilphi}{\tilde{\varphi}}
\newcommand{\opn}{\operatorname}
\newcommand{\re}{\operatorname{Re}}
\newcommand{\im}{\operatorname{Im}}
\newcommand{\diag}{\operatorname{diag}}
\newcommand{\phm}{\phantom{-}}
\newcommand{\mbb}[1]{\mathbb{#1}}
\newcommand{\mb}[1]{\mathbf{#1}}
\newcommand{\mc}[1]{\mathcal{#1}}
\newcommand{\argsr}{(\theta,\partial_\theta,r\partial_r)}
\newcommand{\argsl}{(\theta,\partial_\theta,\lambda)}
\newcommand{\rdr}{r\partial_r}
\newcommand{\jd}{\displaystyle}
\newcommand{\js}{\scriptstyle}
\newcommand{\jt}{\textstyle}
\newcommand{\jss}{\scriptscriptstyle}
\newcommand{\lami}{(\lambda-\lambda_i)}
\newcommand{\lamo}{(\lambda-\lambda_0)}
\newcommand{\der}[2]{\frac{\partial #1}{\partial #2}}
\newcommand{\derr}[2]{\frac{\delta #1}{\delta #2}}
\newcommand{\half}{\frac{1}{2}}
\newcommand{\lla}{\left\langle}
\newcommand{\rra}{\right\rangle}
%\newcommand{\blue}[1]{\textcolor{blue}{#1}}
\newcommand{\pa}{\partial}
\newcommand{\tr}{\operatorname{tr}}
\newcommand{\la}{\langle}
\newcommand{\ra}{\rangle}

%\newcommand{\x}{\tilde{x}}

\newcommand{\y}{\tilde{y}}
\newcommand{\z}{\tilde{z}}
\renewcommand{\L}{\tilde{L}}
\newcommand{\n}{\mathbf{\tilde{n}}}
\newcommand{\PP}{\operatorname{P}}
\newcommand{\PR}{\operatorname{R}}
\DeclareMathOperator*{\argmax}{argmax}
\newcommand{\bds}{\boldsymbol}
\newcommand{\emach}{\epsilon_{\text{machine}}}
\newcommand{\tol}{\text{tol}}
\newcommand{\leray}{{\text{Leray}}}
\newcommand{\TG}{{\text{TG}}}
\newcommand{\rand}{{\text{rand}}}
\newcommand{\low}{{\text{low}}}
\newcommand{\up}{{\text{up}}}
\newcommand{\Kerr}{{\text{K}}}
\newcommand{\filter}{{\text{filter}}}
\newcommand{\Hou}{{\text{H}}}
\newcommand{\Trel}{{T_{\text{rel}}}}

\def\bnabla{\boldsymbol{\nabla}}
\def\u{\boldsymbol{u}}
\def\vv{\boldsymbol{v}}
\def\w{\boldsymbol{w}}
\def\d{\boldsymbol{d}}
\def\j{\boldsymbol{j}}
\def\bk{\boldsymbol{k}}
\def\n{\boldsymbol{n}}
\def\x{\boldsymbol{x}}
\def\z{\boldsymbol{z}}
\def\bbeta{\boldsymbol{\eta}}
\def\bxi{\boldsymbol{\xi}}
\def\bomega{\boldsymbol{\omega}}
\def\RR{\mathbb{R}}
\def\TT{\mathbb{T}}
\def\CC{\mathbb{C}}

\def\tPhi{\widetilde{\Phi}}
\def\teta{\widetilde{\boldsymbol{\eta}}}

\def\M{\mathcal{M}}
\def\T{\mathcal{T}}
\newcommand{\xinyu}[1]{{\color{blue} (\sf $\clubsuit$ Xinyu: #1)}}
\newcommand{\bartosz}[1]{{\color{red} (\sf $\diamondsuit$ Bartosz: #1)}}

\newcommand{\revt}[1]{{\color{blue}#1}}
%\newcommand{\revt}[1]{{\color{black}#1}}
\newcommand{\revtt}[1]{{\color{red}#1}}
\newcommand{\revtbn}[1]{{\color{magenta}#1}}
%-------------------------------------------------------------------------------

\begin{document}
\author[]{}
\address{}
\email{}
%-------------------------------------------------------------------------------
\keywords{}
\subjclass[]{}
%-------------------------------------------------------------------------------
\title{Euler-Voigt Equations}
\maketitle
%-------------------------------------------------------------------------------

%-------------------------------------------------------------------------------
\section{Introduction}
\label{sec:intro}
%-------------------------------------------------------------------------------

We study the Euler-Voigt equations posed on a periodic domain $\mbb T^3 := (\mbb R/\mbb Z)^3 $
%
\begin{equation}\label{eq:Euler:Voigt}
\begin{aligned}
-\alpha^2 \partial_t\Delta\bds u_\alpha + \partial_t \bds u_\alpha + (\bds u\cdot\nabla)\bds u_\alpha + \nabla p &= 0, \\
\nabla \cdot \bds u_\alpha &= 0, \\
\bds u_\alpha(\bds x, 0) &= \bds \eta_\alpha,
\end{aligned}
\end{equation}
%
where $\alpha > 0$ is a regularization parameter.
The Euler-Voigt equations can be viewed as an inviscid regularization of the Euler equations
%
\begin{equation}\label{eq:Euler}
\begin{aligned}
\partial_t \bds u + (\bds u\cdot\nabla)\bds u + \nabla p &= 0, \\
\nabla \cdot \bds u &= 0, \\
\bds u(\bds x, 0) &= \bds \eta.
\end{aligned}
\end{equation}
%
The solutions to \eqref{eq:Euler} and \eqref{eq:Euler:Voigt} both possess the following time scaling symmetry
%
\begin{equation}\label{eq:sym}
\tilde {\bds u}(\bds x, t) = \lambda\bds u(\bds x, \lambda t).
\end{equation}
%

Unlike the Euler equations for which the global well-posedness of the solutions remains an open problem, the solutions of the Euler-Voigt are known to be globally well-posed for $t \in (-\infty, \infty)$ \cite{CaoLunasinTiti2006}. Moreover, the ``$\alpha$-energy'' of the solutions of  \eqref{eq:Euler:Voigt} is conserved for all times
%
\begin{equation}
E_\alpha(t) := ||\bds u_\alpha(t)||^2_{L^2} + \alpha^2 ||\nabla \bds u_\alpha(t)||_{L^2}^2 = E_\alpha(0). 
\end{equation}
%

The behavior of the solutions of \eqref{eq:Euler:Voigt} as $\alpha \to 0$ can be used to determine whether the solutions of the Euler equations will develop a finite time singularity \cite{Larios2018}. 
%
\begin{thm}
Let $\bds \eta_\alpha \in H^s, s \geq 3$ with $\nabla \cdot \bds \eta_\alpha  = 0$ and let $\bds u_\alpha$ be the corresponding unique solution of \eqref{eq:Euler:Voigt}. Suppose that there exits $T^* > 0$ such that
%
\begin{equation}
\limsup_{\alpha\to 0^+} \left(\alpha \sup_{t\in[0, T^*]} ||\nabla \bds u_\alpha(t)||_{L^2}\right) > 0.
\end{equation}
%
Then the Euler equations with the initial condition $\bds \eta_\alpha$ must develop a singularity within the interval $[0, T^*]$.
\end{thm}
%

The Theorem 5.1 in \cite{Larios:Titi:2010} stated that
%
\begin{thm} \label{thm:convergence}
Given initial data $\bbeta$, $\bbeta_\alpha \in H^s(\mbb T^3)$ for some $s\geq 3$, let $\bds u, \bds u_\alpha \in C([0, T], H^s(\mbb T^3)) \cap C^1([0, T], H^{s-1}(\mbb T^3))$ be the corresponding solutions
to \eqref{eq:Euler} and \eqref{eq:Euler:Voigt}, respectively, where $0 < T < T_{\max}$ and $T_{\max}$ is the maximal time for which a solution to \eqref{eq:Euler} exists and is unique. For all $t\in[0, T]$, we have
%
\begin{equation}
\begin{aligned}
&\quad ||\bds u(t) - \bds u_\alpha(t)||_{L^2}^2 
+ \alpha^2 ||\nabla(\bds u(t)- \bds u_\alpha(t))||_{L^2}^2 \\
&\leq \left(||\bds \eta - \bds \eta_\alpha||_{L^2}^2 
+ \alpha^2 ||\nabla(\bds \eta - \bds \eta_\alpha)||_{L^2}^2\right)e^{Ct}
+ C\alpha^2(e^{Ct} - 1).
\end{aligned}
\end{equation}
%
\end{thm}
Based on Theorem \ref{thm:convergence}, we know that if 
%
\begin{equation}
||\bds \eta - \bds \eta_\alpha||_{L^2}^2 
+ \alpha^2 ||\nabla(\bds \eta - \bds \eta_\alpha)||_{L^2}^2 \to 0, \quad \text{as } \alpha \to 0, 
\end{equation}
%
then for all $t\in[0, T]$, 
%
\begin{equation}
||\bds u(t) - \bds u_\alpha(t)||_{L^2}^2 
+ \alpha^2 ||\nabla(\bds u(t) - \bds u_\alpha(t))||_{L^2}^2 \to 0.
\end{equation}
%
Therefore, we can put the constraint on $||\nabla\bbeta_\alpha||_{L^2}$ in the optimization process and check whether 
$||\bds \eta - \bds \eta_\alpha||_{L^2}^2 \to 0$ as $\alpha \to 0$ in computation.

We slightly modify Theorem 5.2 in \cite{Larios:Titi:2010} as follows.
%
\begin{thm}\label{thm:blowup}
Given initial data $\bbeta$, $\bbeta_\alpha \in H^s(\mbb T^3)$ for some $s\geq 3$ satisfying
%
\begin{equation}
||\bds \eta - \bds \eta_\alpha||_{L^2}^2 
+ \alpha^2 ||\nabla(\bds \eta - \bds \eta_\alpha)||_{L^2}^2 \to 0, \quad \text{as } \alpha \to 0,
\end{equation}
%
let $\bds u, \bds u_\alpha \in C([0, T], H^s(\mbb T^3)) \cap C^1([0, T], H^{s-1}(\mbb T^3))$ be the corresponding solutions to \eqref{eq:Euler} and \eqref{eq:Euler:Voigt}, respectively. 

If there exists a finite time $T^* > 0$ such that the solutions $\bds u_\alpha$ satisfies
%
\begin{equation}
\limsup_{\alpha\to 0} \left(\alpha^2 \sup_{t \in [0, T^*]}||\nabla \bds u_\alpha(t)||^2_{L^2}\right)> 0,
\end{equation}
%
then $\bds u$ develops a singularity in the interval $[0, T^{*}]$.
\end{thm}
%

\begin{proof}
Suppose that the solution $\bds u$ of the Euler equations \eqref{eq:Euler} with the initial condition $\bbeta$ is regular
on $[0, T^*]$.
Since the Euler-Voigt equations are globally well-posed, we have 
%
\begin{equation}
||\bds u_\alpha(t)||^2_{L^2} + \alpha^2 ||\nabla \bds u_\alpha(t)||^2_{L^2}  = 
||\bbeta_\alpha||^2_{L^2} + \alpha^2 ||\nabla\bbeta_\alpha||^2_{L^2},
\end{equation}
%
thus
%
\begin{equation}
 \alpha^2 ||\nabla \bds u_\alpha(t)||^2_{L^2}  = 
||\bbeta_\alpha||^2_{L^2} + \alpha^2 ||\nabla\bbeta_\alpha||^2_{L^2} - ||\bds u_\alpha(t)||^2_{L^2}.
\end{equation}
%
Taking the supremum over $t\in[0,T^*]$ of the above equation, we obtain that
 %
\begin{equation}
\sup_{t\in[0,T^*]}\alpha^2 ||\nabla \bds u_\alpha(t)||^2_{L^2} = 
||\bbeta_\alpha||^2_{L^2} + \alpha^2 ||\nabla\bbeta_\alpha||^2_{L^2} - \inf_{t\in[0, T^*]} ||\bds u_\alpha(t)||^2_{L^2}.
\end{equation}
%
Taking the $\limsup$ when $\alpha \to 0$ of the above equation, we obtain that
%
\begin{equation}
\begin{aligned}
\limsup_{\alpha\to 0} \sup_{t\in[0,T^*]}\alpha^2 ||\nabla \bds u_\alpha(t)||^2_{L^2} 
&= \limsup_{\alpha\to 0} \left(||\bbeta_\alpha||^2_{L^2} + \alpha^2 ||\nabla\bbeta_\alpha||^2_{L^2} - \inf_{t\in[0, T^*]} ||\bds u_\alpha(t)||^2_{L^2}\right)\\
&= ||\bbeta||^2_{L^2} - ||\bbeta||^2_{L^2} \\
& = 0,
\end{aligned}
\end{equation}
%
which contradicts the hypothesis in the theorem that the above expression is positive.
The second equality holds because
%
\begin{align*}
\limsup_{\alpha\to 0} \left|\inf_{t\in[0, T^*]} ||\bds u_\alpha(t)||_{L^2} - \| \bbeta\|_{L^2}\right| 
\leq  &\limsup_{\alpha\to 0}||\bds u_\alpha(t_\alpha) - \bbeta||^2_{L^2}  \\
\leq & \limsup_{\alpha\to 0} ||\bds u_\alpha(t_\alpha) - \bds u(t_\alpha)||^2_{L^2}
+ ||\bds u(t_\alpha) - \bbeta||^2_{L^2} \\
= &0.
\end{align*}
%
Here in the second inequality, we use the fact that the convergence of the solutions of the Euler-Voigt equations to the solution of the Euler equations is uniform on $[0, T^*]$, as shown in Theorem \ref{thm:convergence}.




\end{proof}
 
 
%-------------------------------------------------------------------------------
\section{PDE optimization method}
%-------------------------------------------------------------------------------

%-------------------------------------------------------------------------------
\subsection{The optimization problem}
%-------------------------------------------------------------------------------


We consider analytic initial conditions that belong to the following Gevrey class
\begin{equation}
G^\sigma := \left\{ \vv \in H^s(\TT^3) \ : \ \| \vv \|_{G^\sigma} = \left\langle \vv,\vv \right\rangle^{1/2}_{G^\sigma} < \infty \right\}, \quad s \geq 3, \quad \sigma > 0,
\label{eq:G}
\end{equation}
with the inner product
\begin{equation}\label{eq:Gevrey:def}
\begin{aligned}
\forall \bds v, \bds u\in G^{\sigma}, \quad 
\left\la\bds v, \bds u \right\ra_{G^{\sigma}} 
:= &\sum_{\bds j \in \mathbb Z^3} 
(1+|2\pi\bds j|^2)^s e^{4\pi\sigma|\bds j|} \hat{\bds v}_{\bds j} \cdot \overline{\hat{\bds u}}_{\bds j}\\
=& \int_{\mbb T^3}(1+|D|^2)^s e^{2\sigma|D|} \bds v \cdot \bds u \, d\bds x,
\end{aligned}
\end{equation}
where overbar denotes complex conjugation, whereas the operators $|D|$
and $e^{\sigma|D|}$ are defined via
%
\begin{equation}
\left[\widehat{|D| \bds v}\right]_{\bds j} := 2\pi|\bds j| \hat{\bds v}_{\bds j}, \qquad \qquad 
\left[\widehat{e^{\sigma |D|}\bds v}\right]_{\bds j} := e^{2\pi\sigma|\bds j|}\hat{\bds v}_{\bds j}.
\label{eq:D}
\end{equation}
Since the initial conditions in \eqref{eq:Euler:Voigt} and \eqref{eq:Euler} need to
be divergence-free and the quantity $\int_{\TT} \bbeta \, d\x$ is an
invariant of motion, we introduce the following subspace in which
optimal initial conditions will be sought
%
\begin{equation}\label{eq:space:initial:condition}
V := \{ \bds v  \in G^{\sigma}: \,
\bnabla \cdot \bds v = 0, \quad \int_{\mbb T^3} \bds v \, d \bds x= \bds 0\}.
\end{equation}
%
According to the analytical results in Section \ref{sec:intro}, we define the objective functional $\Phi_{\alpha, T}: V \to \mbb R^+$ as
%
\begin{equation}\label{eq:obj}
\Phi_{\alpha, T}(\bbeta) := \| \nabla\u_\alpha(T; \bbeta) \|^2_{L^2} 
= \| \u_\alpha(T; \bbeta) \|^2_{\dot{H}^1},
\end{equation}
%
where $\u_\alpha(T;\bbeta)$ is the solution of the Euler-Voigt equations
\eqref{eq:Euler:Voigt} obtained with the initial condition $\bbeta \in V$ and $\alpha$ is the regularization parameter in \eqref{eq:Euler:Voigt}. 

Due to the scaling property \eqref{eq:sym} of the Euler-Voigt equations \eqref{eq:Euler:Voigt},  
we can restrict our discussion to initial conditions with
unit $\dot H^1$ seminorm which belong to a closed manifold
$\M_1 \subset V$ defined as
\begin{equation}\label{eq:M1}
\M_1 = \{\bds v \in V:\,  ||\bds v||_{\dot H^1} = 1\}.
\end{equation}
Therefore, we  arrive at the following optimization problem
%
\begin{problem}\label{pb:H3}
Given $\alpha, T \in \mbb R_+$, find 
%
\begin{equation}
\teta_{\alpha, T} = \argmax\limits_{\bbeta \in \M_{1}} \Phi_{\alpha, T}(\bbeta).
\end{equation}
%
\end{problem}
Our goal is to first fix $T$ and solve the optimization problem \ref{pb:H3} for a sequence of decreasing values of $\alpha$ that converge to 0.


%-------------------------------------------------------------------------------
\subsection{The adjoint method}
%-------------------------------------------------------------------------------
We compute the gradient $\nabla \Phi_{\alpha, T}(\bbeta)$ of the objective
functional in \eqref{eq:obj} using an adjoint method which depends on solving a properly defined
adjoint system backward in time. We 
first introduce the G\^{a}teaux
(directional) differential $\Phi_{\alpha, T}'(\bbeta; \bbeta'): V \times V
\to \mbb R$
%
\begin{equation}\label{eq:differential}
\begin{aligned}
\Phi_{\alpha, T}'(\bbeta; \bbeta') :&= \lim_{\epsilon \to 0} \frac{1}{\epsilon}
\left[ \Phi_{\alpha, T}\left(\bbeta + \epsilon \bbeta'\right) - \Phi_{\alpha, T}(\bbeta)\right]\\
& = 2\left\la\bds u_\alpha(\cdot, T), \bds u_\alpha'(\cdot, T)\right\ra_{\dot H^1}.
\end{aligned}
\end{equation}
Here $\bds u_\alpha'(\bds x, t)$ is the solution of the linearization of the
Euler-Voigt equations \eqref{eq:Euler:Voigt} around its solution $\bds u_\alpha(\bds x, t;
\bbeta)$, which is defined by
%
\begin{equation}\label{eq:L}
\begin{aligned}
\mc L
\begin{bmatrix} \bds{u}_\alpha' \\ p_\alpha' \end{bmatrix}
:&=
\begin{bmatrix}
\partial_t \bds u_\alpha' + (I-\alpha^2\Delta)^{-1}\left(\bds u_\alpha' \cdot \bnabla \bds u_\alpha + \bds u_\alpha\cdot \bnabla \bds u_\alpha' + \bnabla p_\alpha'\right)\\[3pt]
\bnabla \cdot \bds u_\alpha'
\end{bmatrix}=
\begin{bmatrix}
\bds 0\\[3pt]
0
\end{bmatrix},\\[3pt]
\u_\alpha'(\bds x, 0) & = \bbeta'(\bds x),
\end{aligned}
\end{equation}
%
where $\bbeta'$ is the perturbation of the initial condition and
$p'_\alpha(\bds x, t)$ is the corresponding pressure perturbation.
We further introduce the adjoint states $\left\{\bds u^*_\alpha: \mbb T^3 \times [0, T]
  \to \mbb R^3, \ p^*_\alpha: \mbb T^3 \times [0, T] \to \mbb R\right\}$ and
the following duality-pairing relation
%
\begin{equation}\label{eq:pairing}
\begin{aligned}
\left(\mc L 
\begin{bmatrix} \bds{u}_\alpha' \\ p_\alpha' \end{bmatrix}, \, 
\begin{bmatrix} \bds{u}_\alpha^* \\ p_\alpha^* \end{bmatrix}
\right) : =& 
\int_0^{T}\int_{\mbb T^3} 
\mc L \begin{bmatrix} \bds{u}_\alpha' \\ p_\alpha' \end{bmatrix} \cdot 
\begin{bmatrix} \bds{u}_\alpha^* \\ p_\alpha^* \end{bmatrix}
\, d\bds x \,dt\\[3pt]
= &\left(
\begin{bmatrix} \bds{u}_\alpha' \\ p_\alpha' \end{bmatrix}, \, 
\mc L^* \begin{bmatrix} \bds{u}_\alpha^* \\ p_\alpha^* \end{bmatrix}
\right)
+ \int_{\mbb T^3} \bds u_\alpha'(\bds x, T) \cdot \bds u_\alpha^*(\bds x, T) \, d\bds x\\[3pt]
&- \int_{\mbb T^3} \bbeta'(\bds x) \cdot \bds u_\alpha^*(\bds x, 0) \, d\bds x\\[3pt]
= &\  0.
\end{aligned}
\end{equation}
%
Here ``$\cdot$'' denotes the usual Euclidean inner product on
$\RR^4$ and $\mc L^*$ is the {\em adjoint} operator defined
in terms of the following system, which is obtained by performing
integration by parts with respect to both space and time in
(\ref{eq:pairing})
%
\begin{equation}\label{eq:adjoint}
\begin{aligned}
\mc L^* \begin{bmatrix} \bds{u}_\alpha^* \\ p_\alpha^* \end{bmatrix} :&= 
\begin{bmatrix}
-\partial_t \bds u_\alpha^* - (\bnabla (I - \alpha^2\Delta)^{-1}\bds u_\alpha^* + 
(\bnabla(I-\alpha^2\Delta)^{-1}\bds u_\alpha^*)^T)\bds u_\alpha - \bnabla p_\alpha^*\\[3pt]
-\bnabla \cdot \bds u_\alpha^*
\end{bmatrix}
= \begin{bmatrix}
\bds 0\\[3pt]
0
\end{bmatrix}, \\[3pt]
\bds u^*_\alpha(\bds x, T) &=  2|D|^2\bds u_\alpha(\bds x, T).
\end{aligned} 
\end{equation}
%
Therefore, we have
%
\begin{equation}
\Phi_{\alpha, T}'(\bbeta; \bbeta') = 2\left\la\bds u_\alpha(\cdot, T), \bds u_\alpha'(\cdot, T)\right\ra_{\dot H^1} = \left\la\bds u_\alpha^*(\cdot, 0), \bds \eta'(\cdot)\right\ra_{L^2} 
\end{equation}
%







%%%%%%%%%%%%%%%%%%%%%%%%%%%%%%%%%%%%%%%%%

\bibliographystyle{plain}
\bibliography{references} 


\end{document}

%% To be filled in the journal office:

@author:
@affiliation:
@title:
@language: English
@pages:
@classification1:
@classification2:
@keywords:
@abstract:
@filename:
@EOI


